\documentclass[aspectratio=169,onlytextwidth]{beamer}

\usetheme[english]{uzh} % Institution label: english or german

% Additional packages for the demo content
\usepackage{amsmath,amssymb}
\usepackage{booktabs}
\usepackage{uzhcode} % UZH-styled source-code listings + \cmd / \env helpers
\usepackage{uzhdag}  % TikZ styles for Bayesian-network / DAG diagrams

% In-text \cite marks become small, raised and grey -- pointers, not content.
% The reference list keeps its normal size. Remove the line for plain [1].
\uzhtinycites

\title{A Beamer Theme for the University of Zurich}
\subtitle{THU-style features in the UZH corporate design}
\author{Hans Muster}
\date{\today}
% The institute goes in the title-page header (it wraps by itself -- do not
% put \\ in it). The footline shows the UZH / DM3L wordmark instead of the
% university and institute names, so \institute no longer appears there.
\institute{Department of XYZ}
\occasion{1st Committee Meeting} % purpose of the talk, shown in the footline
\titlegraphic{img/math_institute.jpg}

\begin{document}

\maketitle

% Overview slide: sections only (see \uzhoutlineframe in uzhnavigation.sty)
\uzhoutlineframe


\section{Introduction}

\subsection{About this theme}

\begin{frame}{What this theme gives you}
  This template combines the \textbf{look} of the official UZH corporate-design
  beamer theme with the \textbf{navigation features} of the popular THU beamer
  theme.

  \begin{itemize}
    \item A section navigation bar with progress dots at the top of every slide
    \item Automatic \emph{outline} slides at the start of each section
    \item A footer showing the current slide and the total number of slides
    \item Coloured blocks, code listings and math, all in the UZH palette
  \end{itemize}

  Everything is built on the standard \texttt{beamer} class~\cite{lamport1994latex},
  so all usual commands and environments keep working.
\end{frame}

\begin{frame}{\steptitle{Building up a list step by step}}
  \setsteps{3}
  \stepnote{Reveal one point at a time with \texttt{steplist}: the current
    point is black, earlier points fade to grey, and a final overlay shows the
    whole slide in black to talk over as a summary.}

  \begin{steplist}
    \stepitem{First the motivation becomes visible \ldots}
    \stepitem{\ldots then the method \ldots}
    \stepitem{\ldots and finally the result.}
  \end{steplist}

  \stepsummary{Use \texttt{\textbackslash pause} or \texttt{[<+->]} for a plain
    reveal.}
\end{frame}

\begin{frame}{\steptitle{Stepping through blocks and prose}}
  % \setsteps{1} puts the closing all-black overlay at 1+3 = 4, which gives the
  % three content overlays this frame uses.
  \setsteps{1}
  \begin{columns}[T]
    \begin{column}{0.47\textwidth}
      \uncover<2->{{\stepblockfade{2}%
      \begin{block}{\texttt{\textbackslash stepblockfade}}
        Fades a whole block. It is a box, so reveal it with
        \texttt{\textbackslash uncover}.
      \end{block}}}
    \end{column}
    \begin{column}{0.47\textwidth}
      \onslide<3->%
      {\steptextfade{3}%
      \texttt{\textbackslash steptextfade} fades plain prose, and display math
      goes with it:
      \[ a^2 + b^2 = c^2 \]
      Reveal prose with \texttt{\textbackslash onslide}, because
      \texttt{\textbackslash uncover} breaks display math.}
    \end{column}
  \end{columns}

  \stepsummary{Both helpers only set the colour; you choose the overlay.}
\end{frame}

\subsection{Colours}

\begin{frame}{The UZH colour palette}
  The corporate colours are available as commands and as named colours:

  \begin{itemize}
    \item \uzhblue{Blue} \quad \uzhcyan{Cyan} \quad \uzhapple{Apple}
    \item \uzhgold{Gold} \quad \uzhorange{Orange} \quad \uzhberry{Berry}
  \end{itemize}

  Use \texttt{\textbackslash textcolor\{uzh@blue\}\{\ldots\}} or the shortcuts
  \texttt{\textbackslash uzhblue\{\ldots\}} shown above. Shades 1--5 exist for
  each colour (e.g. \texttt{uzh@blue1} \ldots\ \texttt{uzh@blue5}).
\end{frame}


\section{Content elements}

\subsection{Blocks}

\begin{frame}{Blocks}
  \begin{block}{Definition}
    A \emph{block} groups related content under a coloured title bar.
  \end{block}

  \begin{exampleblock}{Example}
    Example blocks use the UZH apple colour.
  \end{exampleblock}

  \begin{alertblock}{Warning}
    Alert blocks use the UZH berry colour to draw attention.
  \end{alertblock}
\end{frame}

\subsection{Lists, columns and images}

\begin{frame}{Text and a picture}
  \begin{columns}
    \begin{column}{0.62\textwidth}
      Combine running text with images using \texttt{columns}. The
      \texttt{\textbackslash fillimage} helper crops an image to exactly fill a
      given width and height, so photos never look stretched.

      \begin{enumerate}
        \item Pick a width and a height
        \item The image is centred and cropped to fit
        \item Add an \texttt{alt} text for accessibility
      \end{enumerate}
    \end{column}
    \begin{column}{0.34\textwidth}
      \fillimage[alt={A staircase at UZH}]{\textwidth}{6cm}{img/uzh-staircase.jpg}
    \end{column}
  \end{columns}
\end{frame}

\subsection{DAG diagrams}

\begin{frame}{Bayesian-network diagrams}
  Load \texttt{uzhdag} for TikZ styles in the UZH palette: \texttt{dagnode},
  \texttt{dagedge} and \texttt{dagabsent} (dashed berry $=$ absent edge).

  \begin{columns}
    \begin{column}{0.42\textwidth}
      \centering
      \begin{tikzpicture}[x=2.1cm, y=1.6cm]
        \node[dagnode] (am) at (0,0)  {$a_{t-1}$};
        \node[dagnode] (bm) at (0,-1) {$b_{t-1}$};
        \node[dagnode] (a0) at (1,0)  {$a_{t}$};
        \node[dagnode] (b0) at (1,-1) {$b_{t}$};
        \draw[dagedge]   (am) -- (a0);
        \draw[dagedge]   (bm) -- (b0);
        \draw[dagedge, line width=2pt] (am) -- (bm);
        \draw[dagabsent] (am) -- (b0);
        \draw[dagabsent] (bm) -- (a0);
      \end{tikzpicture}
    \end{column}
    \begin{column}{0.52\textwidth}
      \begin{itemize}
        \item solid arrow: a present edge
        \item dashed berry arrow: an absent edge --- ``missing arrows are
          the information''
        \item extra TikZ options compose, e.g.\
          \texttt{line width=2pt} for a strong coefficient
      \end{itemize}
    \end{column}
  \end{columns}
\end{frame}

\begin{frame}{Profile cards}
  \begin{columns}
    \begin{column}[T]{0.23\textwidth}
      \imagecard{img/profile.pdf}{Martha Muster}{Professor}
    \end{column}
    \begin{column}[T]{0.23\textwidth}
      \imagecard{img/profile.pdf}{Hans Muster}{Academic Associate}
    \end{column}
    \begin{column}[T]{0.23\textwidth}
      \imagecard{img/profile.pdf}{Martha Muster}{PhD Student}
    \end{column}
    \begin{column}[T]{0.23\textwidth}
      \imagecard{img/profile.pdf}{Hans Muster}{PhD Student}
    \end{column}
  \end{columns}
\end{frame}


\section{Typesetting}

\subsection{Mathematics}

\begin{frame}{Mathematics}
  Inline math such as $\int_0^1 x^2 \, dx = \tfrac{1}{3}$ flows with the text.
  Display math is centred:
  \[
    \sum_{i=1}^{n} i = \frac{n(n+1)}{2}.
  \]

  Aligned, numbered equations work as usual:
  \begin{align}
    Q_{\text{target}} &= r + \gamma\, Q^\pi\big(s', \pi_\theta(s') + \epsilon\big) \\
    \epsilon &\sim \mathrm{clip}\big(\mathcal{N}(0,\sigma), -c, c\big) \nonumber
  \end{align}
\end{frame}

\subsection{Tables and code}

\begin{frame}{Tables}
  Use \texttt{booktabs} for clean, professional tables:

  \begin{table}
    \centering
    \begin{tabular}{lcc}
      \toprule
      Method & Accuracy & Runtime (ms) \\
      \midrule
      Baseline      & 0.812 & 14.2 \\
      \textbf{Ours} & \textbf{0.947} & \textbf{9.8} \\
      Prior work    & 0.903 & 21.6 \\
      \bottomrule
    \end{tabular}
    \caption{Comparison of methods on the benchmark.}
  \end{table}
\end{frame}

\begin{frame}[fragile]{Source code}
  Mark the frame \texttt{[fragile]} and use \texttt{lstlisting} (styled by
  \env{uzhcode}) to show code:

\begin{lstlisting}[language=Python]
def fib(n):
    # return the n-th Fibonacci number
    a, b = 0, 1
    for _ in range(n):
        a, b = b, a + b
    return a
\end{lstlisting}

  Inline, \cmd{section} starts a section and \env{itemize} an unnumbered list.
\end{frame}


\section{Full-bleed slides}

\begin{frame}[t]
  \numberoverlay{42}
\end{frame}

\begin{frame}
  \vskip0.4cm
  \fillimage{\textwidth}{\textheight}{img/uzh-city.jpg}
\end{frame}


\section{References}

\nocite{*} % list every entry in ref.bib, even the uncited ones

% The reference list on one slide, in two columns. Write \uzhbibframe[1]{ref}
% for a single column, or a plain frame with [allowframebreaks] if the list is
% too long for one slide.
\uzhbibframe{ref}

% The closing page. Rename either line with \renewcommand{\uzhthanksname}{...}
% or \renewcommand{\uzhquestionsname}{...} in the preamble.
\uzhthanksframe


% ===========================================================================
% Appendix -- pages that are kept for a question and are not walked through.
% \appendix opens a new part, thus these frames stay out of the table of
% contents of the opening overview. \section* names the part in the navigation
% bar without adding a contents entry.
%
% The footline counts these frames on their own ("3 / 4" and not "34 / 37"),
% and the navigation bar carries a button that jumps back to the opening
% overview -- from there every entry is a link into the talk.
% ===========================================================================
\appendix
\section*{Appendix}

\begin{frame}{A page kept for a question}
  A frame after \texttt{\textbackslash appendix} is shown only if somebody
  asks for it.

  \begin{itemize}
    \item The footline counts the appendix frames from 1 again
    \item The navigation bar shows a \emph{Contents} button at its right end
    \item The button jumps to the opening overview, and every entry there is a
      link --- thus two clicks reach any part of the talk
  \end{itemize}

  Rename the button with
  \texttt{\textbackslash renewcommand\{\textbackslash uzhtocbuttonname\}\{Overview\}}.
\end{frame}

\end{document}
